\section{The decoded skeleton and its minor ranges}
\label{sec:decoded-skeleton}

After \cref{prop:global-fibre-error}, it remains to estimate the single decoded tuple attached to each anchor assignment.  Coherent anchor assignments are indexed by an integer label $m$; noncoherent assignments are treated by the anchor partition.  The lower--lower pair factors were retained during fibre compression.

\subsection{Prime-only Gaussian transition}

Assume $1<\gamma\leq2$ and
\begin{equation}
\label{eq:transition-range}
 T_{\mathrm{full}}<|m|\leq X^2/4.
\end{equation}
By \cref{prop:prime-decoder-identification}, every coordinate indexed by $P$ equals $m$.  Hence every complete-pair factor has phase $m/(pq)$.  Since $pq\geq X^2$ and $|m|\leq X^2/4$, these phases lie in the linear circle range, and \cref{eq:bernoulli-modulus} gives
\begin{equation}
\label{eq:transition-pointwise}
 |F(h_m)|
 \leq\exp\left(-c m^2
   \sum_{p<q\in P}\frac1{p^2q^2}\right)
 \leq\exp\left(-c\frac{m^2}{X^2\log^2X}\right).
\end{equation}
At the lower endpoint, the exponent satisfies
\begin{equation}
\label{eq:transition-lower-exponent}
 \frac{T_{\mathrm{full}}^2}{X^2\log^2X}
 \gg\frac{Z^2}{X^2\log^2X}
 =\frac{X^{2\gamma-2}}{\log^2X}\longrightarrow\infty.
\end{equation}
Therefore
\begin{equation}
\label{eq:transition-total}
 \sum_{T_{\mathrm{full}}<|m|\leq X^2/4}|F(h_m)|
 =o(\sigma_{\mathcal E}^{-1}).
\end{equation}
This sector uses no target coordinate.  At $\gamma=2$ it treats the variance boundary exactly; for $1<\gamma<2$ it fills the entire gap left by the shorter full-coordinate range.

\subsection{Adaptive retained-pair damping}

For an integer $m$ with
\[
 X^2/4<|m|\leq M_{\mathrm{dec}},
\]
put $I_m=[2\sqrt{|m|},3\sqrt{|m|}]$.

\begin{proposition}[Adaptive prime interval]
\label{prop:adaptive-prime-interval}
For every such $m$, the interval $I_m$ lies in $[X,Z/2)$ and contains
\[
 K_m\gg\frac{\sqrt{|m|}}{\log|m|}
\]
primes.  On the decoded skeleton, the retained complete-pair energy satisfies
\begin{equation}
\label{eq:adaptive-energy}
 Q_{\mathrm{pair}}(m)\gg\frac{|m|}{(\log|m|)^2}.
\end{equation}
\end{proposition}

\begin{proof}
The lower inclusion follows from $2\sqrt{|m|}>X$.  At the upper endpoint,
\[
 3\sqrt{|m|}\leq\frac{3\sqrt{XZ}}{\log Z}<Z/2
\]
for every fixed $\gamma>1$ and large $X$.  Fixed-ratio prime supply gives the count.  For distinct $p,q\in I_m$, one has $1/9\leq|m|/(pq)\leq1/4$, so each retained pair contributes a fixed positive amount of energy; summing over the pairs proves \cref{eq:adaptive-energy}.
\end{proof}

Consequently,
\begin{equation}
\label{eq:adaptive-pointwise}
 |F(h_m)|\leq\exp\left(-c\frac{|m|}{(\log|m|)^2}\right),
\end{equation}
and
\begin{equation}
\label{eq:adaptive-total}
 \sum_{X^2/4<|m|\leq M_{\mathrm{dec}}}|F(h_m)|
 =o(\sigma_{\mathcal E}^{-1}).
\end{equation}

\subsection{Pairwise-disjoint exhaustive partitions}

Fix $C\geq1$, to be chosen before $X$, and put
\begin{equation}
\label{eq:N-major}
 N=\left\lfloor\frac{C}{\sigma_{\mathcal E}}\right\rfloor.
\end{equation}
By \cref{prop:variance-regimes}, for every fixed $C$,
\begin{equation}
\label{eq:N-below-full}
 N=o(T_{\mathrm{full}}).
\end{equation}

Apply the decoder/nondecoder split first.  Every nondecoder tuple belongs to the terminal error sector.  Among decoded tuples, a noncoherent anchor assignment belongs to that same sector.  A coherent assignment has one unique label $m$.

If $\gamma>2$, the following five sectors are pairwise disjoint and exhaustive:
\begin{center}
\begin{tabular}{@{}c p{5.2cm} p{5.0cm}@{}}
\toprule
Sector & skeleton condition & control \\
\midrule
I & $|m|\leq N$ & positive Taylor major \\
II & $N<|m|\leq X^2/4$ & actual-variance Gaussian tail \\
III & $X^2/4<|m|\leq M_{\mathrm{dec}}$ & adaptive retained pairs \\
IV & $|m|>M_{\mathrm{dec}}$ & internal anchor energy \\
V & noncoherent decoded tuple, or any nondecoder tuple & anchor entropy and fibre error \\
\bottomrule
\end{tabular}
\end{center}

If $1<\gamma\leq2$, the following six sectors are pairwise disjoint and exhaustive:
\begin{center}
\begin{tabular}{@{}c p{5.2cm} p{5.0cm}@{}}
\toprule
Sector & skeleton condition & control \\
\midrule
I & $|m|\leq N$ & positive Taylor major \\
II & $N<|m|\leq T_{\mathrm{full}}$ & actual-variance Gaussian tail \\
III & $T_{\mathrm{full}}<|m|\leq X^2/4$ & prime-only pair Gaussian \\
IV & $X^2/4<|m|\leq M_{\mathrm{dec}}$ & adaptive retained pairs \\
V & $|m|>M_{\mathrm{dec}}$ & internal anchor energy \\
VI & noncoherent decoded tuple, or any nondecoder tuple & anchor entropy and fibre error \\
\bottomrule
\end{tabular}
\end{center}

\begin{proposition}[Regime-dependent sector exhaustion]
\label{prop:five-sector-exhaustion}
The five sectors for $\gamma>2$, and the six sectors for $1<\gamma\leq2$, are respectively pairwise disjoint and exhaustive over all frequencies modulo $L$.
\end{proposition}

\begin{proof}
The CRT coordinate map is a bijection.  The decoder/nondecoder split is disjoint for each anchor assignment.  A decoded coherent assignment has exactly one label and therefore lies in exactly one of the displayed disjoint label ranges.  No residual sibling sector remains.
\end{proof}

\subsection{Outer totals}

For $|m|>M_{\mathrm{dec}}$, the internal anchor factor gives
\[
 |F(h_m)|\leq\exp(-\kappa_bm^2\sigma_{B,0}^2).
\]
Thus
\begin{equation}
\label{eq:anchor-tail-total}
 \sum_{|m|>M_{\mathrm{dec}}}|F(h_m)|
 \ll\sigma_{B,0}^{-1}
 \exp(-cM_{\mathrm{dec}}^2\sigma_{B,0}^2)
 =o(\sigma_{\mathcal E}^{-1}),
\end{equation}
because $M_{\mathrm{dec}}^2\sigma_{B,0}^2\gg X^2/\log^6Z$.

Finally, \cref{eq:noncoherent-top-tail,eq:global-fibre-error} give
\begin{equation}
\label{eq:terminal-error-total}
 \sum_{\mathrm{terminal}}|F|
 \ll\exp(-cF_B)+P_{\mathrm{anc}}(\exp(\Delta)-1)
 =o(\sigma_{\mathcal E}^{-1}).
\end{equation}
